% ===== Document Preamble =====
\documentclass[11pt]{article}
\usepackage[margin=0.5in]{geometry}
\usepackage{amsmath,amssymb}
\usepackage{xcolor}
\usepackage{enumitem}
\usepackage{needspace}
\usepackage{comment}

% ===== Colors =====
\definecolor{AnsRed}{RGB}{170,0,0}
\definecolor{WrongBlue}{RGB}{0,0,160}

% ===== Choice list styling (tight, consistent) =====
\setlist[enumerate]{%
  leftmargin=1.5em,
  itemsep=0.25em,
  topsep=0.25em,
  label=\textbf{\Alph*.}
}

% ===== Style Wrappers =====
\newcommand{\QuestionStyle}[1]{{\large\bfseries\boldmath #1}}
\newcommand{\AnswerStyle}[1]{{\large\bfseries\boldmath #1}}
\newcommand{\ReasonStyle}[1]{{\normalsize #1}}

% ===== Show / Hide Reasoning Toggle =====
\newif\ifshowreason
\showreasontrue      % Instructor version
% \showreasonfalse   % Student version

% ===== Choice Item Macros (Answer + Reasoning) =====
% Usage:
%   \Correct{Answer}{Reasoning}
%   \Wrong{Answer}{Reasoning}

\newcommand{\Correct}[2]{%
  \item \AnswerStyle{\textcolor{AnsRed}{#1}}%
  \ifshowreason\\[-0.2em]
    \ReasonStyle{\textcolor{AnsRed}{#2}}%
  \fi
}

\newcommand{\Wrong}[2]{%
  \item \AnswerStyle{\textcolor{WrongBlue}{#1}}%
  \ifshowreason\\[-0.2em]
    \ReasonStyle{\textcolor{WrongBlue}{#2}}%
  \fi
}

% ===== Question Block Macro =====
% Usage: \QBlock{Q1}{stem...}{<items>}

\newcommand{\QBlock}[3]{%
  \Needspace{8\baselineskip}%
  \noindent
  \begin{minipage}{\linewidth}
    \QuestionStyle{#1.~#2}

    \vspace{0.25em}
    \begin{enumerate}
      #3
    \end{enumerate}

    \vspace{0.25em}
    \hrule height 0.4pt
  \end{minipage}
  \vspace{0.75em}
}

% ===== Document Start =====

\begin{document}

\begin{center}
{\huge \textbf{Isomorphic Questions: Vectors and Kinematics}}\\[0.25em]
{\Large Position vectors, displacement, velocity, and acceleration}\\[0.25em]
{\large Multiple choice questions generated below. Answers in RED and incorrect answers in BLUE.}
\end{center}

\vspace{0.75em}

\QBlock{Q1}{A survey drone is launched from a base at the origin. At a certain moment its position vector is $\vec r = 8\hat{i} - 15\hat{j}$, where distances are measured in kilometers. Which statement is correct?}{
  \Wrong{The distance between the drone and the base is $23\ \text{km}$.}{This adds the component magnitudes ($8+15$) as if distance were Manhattan distance. The straight-line distance should be found from the Euclidean magnitude $\sqrt{x^2+y^2}$.}
  \Wrong{The distance between the drone and the base is $15\ \text{km}$.}{This incorrectly uses only the $y$-component and ignores the $x$-component. The distance from the origin depends on both components via $\sqrt{x^2+y^2}$.}
  \Correct{The distance between the drone and the base is $17\ \text{km}$.}{The distance from the origin is the magnitude of the position vector: $|\vec r|=\sqrt{(8)^2+(-15)^2}$. Computing gives $|\vec r|=\sqrt{64+225}=\sqrt{289}=17\ \text{km}$, so this statement is correct.}
  \Wrong{The distance between the drone and the base is $8\ \text{km}$.}{This incorrectly uses only the $x$-component and ignores the $y$-component. The distance from the origin is the vector magnitude, not a single component.}
}

\QBlock{Q2}{A cart moves along the $x$-axis with displacement given by $\vec s(t)=\big((2.2\,\mathrm{m/s^3})t^3+(1.6\,\mathrm{m/s^2})t^2\big)\hat{i}$. Which statement is correct?}{
  \Wrong{This is a constant-acceleration motion.}{Taking two derivatives gives $\vec a(t)=(13.2t+3.2)\hat i$, which depends on time. Since $\vec a$ is not constant, the motion is not constant-acceleration motion.}
  \Correct{The magnitude of the acceleration at $t=0$ is $3.2\,\mathrm{m/s^2}$.}{Differentiate to get velocity $\vec v(t)=\frac{d\vec s}{dt}=(6.6t^2+3.2t)\hat i$ and acceleration $\vec a(t)=\frac{d\vec v}{dt}=(13.2t+3.2)\hat i$. Evaluating at $t=0$ gives $|\vec a(0)|=3.2\,\mathrm{m/s^2}$.}
  \Wrong{The magnitude of the acceleration is $1.6\,\mathrm{m/s^2}$.}{This incorrectly reads a coefficient from the displacement function as the acceleration. Acceleration must be found via $\vec a(t)=\frac{d^2\vec s}{dt^2}$, not by inspection of the $t^2$ term.}
  \Wrong{The magnitude of the acceleration at $t=0$ is $13.2\,\mathrm{m/s^2}$.}{This confuses the coefficient of $t$ in $\vec a(t)=(13.2t+3.2)\hat i$ with the value at $t=0$. Substituting $t=0$ leaves only the constant term $3.2\,\mathrm{m/s^2}$.}
}

\QBlock{Q3}{An object moves in the plane from position $A$ with $\vec r_A=-7.0\hat{i}+4.0\hat{j}$ to position $B$ with $\vec r_B=5.0\hat{i}-2.0\hat{j}$ in $4.0\ \text{s}$. Coordinates are in meters. Which statement is correct?}{
  \Correct{The straight-line distance between $A$ and $B$ is about $13.4\ \text{m}$.}{The displacement is $\Delta\vec r=\vec r_B-\vec r_A=(5-(-7))\hat i+(-2-4)\hat j=12\hat i-6\hat j$. Its magnitude is $|\Delta\vec r|=\sqrt{12^2+(-6)^2}=\sqrt{180}\approx 13.4\,\text{m}$, so this statement is correct.}
  \Wrong{The speed at $t=0$ must be $3.4\ \text{m/s}$.}{This assumes the initial speed must equal the average speed $|\Delta\vec r|/\Delta t$. With only start/end positions and total time, infinitely many motions are possible, so the initial speed is not uniquely determined.}
  \Wrong{The acceleration at $t=0$ must be $0.84\ \text{m/s}^2$.}{This effectively invents an acceleration using a relation like $a\sim |\Delta\vec r|/t^2$ without justification. Acceleration cannot be determined from just two position vectors and an elapsed time unless additional constraints (e.g., constant acceleration plus more information) are provided.}
  \Wrong{None of the above is correct.}{Option A is correct because the magnitude of the displacement between the two given coordinates is about $13.4\,\text{m}$. Therefore it is not true that none of the statements is correct.}
}

\end{document}
